% ============================================================
% VRAM and Iteration-time Scaling — LLOCAMuPTransformer
% Hardware: Tesla V100-SXM2-16GB (Jean-Zay gpu\_p2 partition)
%           for 32 GB projections: V100-SXM2-32GB (gpu\_p2l)
% ============================================================

\section{GPU Scaling Benchmarks}

Jean-Zay allocates GPU hours as a fixed budget per project, so
maximising the amount of useful computation per step is critical.
This section quantifies two bottlenecks: peak VRAM (which limits
batch size and model width) and step time (which limits how many
gradient updates fit in an allocation).  The goal is to understand
the trade-off between batch size $\mathrm{bs}$, number of particles
per event $n$, and network width $n_h$ (number of attention heads,
$d = 16 n_h$), so that each job runs close to the hardware ceiling.

Measurements were collected with \texttt{test\_vram\_scaling.py} and
\texttt{test\_iteration\_time.py}, which perform realistic steps
including a full forward pass, backward pass, and \textsc{AdamW}
update.

\subsection{Peak VRAM}

\subsubsection{Fitting formula}

Peak VRAM during a training step is well described by
\begin{equation}
  V_\text{peak}(\mathrm{bs},\, n,\, n_h)
    \;\approx\; a(n_h)\cdot\mathrm{bs}\cdot n
              + b(n_h)\cdot\mathrm{bs}
              + C(n_h)
  \quad [\text{MB}],
  \label{eq:vram}
\end{equation}
where the $b\cdot\mathrm{bs}$ term is small relative to the leading
term for $n\geq 2$.  The coefficient $a$ accounts for activation
memory (scales as $n_h$) and $C$ for fixed memory: model parameters
plus the two AdamW momentum buffers ($\approx 3\times$ params in
float32).

\subsubsection{Width scaling}

The network width is $d = 16\,n_h$; the framesnet MLP is fixed at
128 hidden channels regardless of width.  Parameters scale roughly
as $n_h^2$ (transformer attention and MLP blocks dominate), so fixed
memory grows as $n_h^2$ while the activation coefficient $a$ grows
linearly with $n_h$.

\begin{table}[h]
\centering
\caption{%
  Model size and VRAM fit coefficients vs.\ width.
  $d = 16\,n_h$ is the transformer hidden dimension.
  Fixed memory = params + AdamW states ($\approx 3\times$ params).
  EMA (\texttt{torch\_ema}) adds one float32 shadow copy on GPU,
  increasing $C$ by exactly $\Delta C_\text{EMA} = \text{params (MB)}$;
  $a$ and $b$ are unchanged.
  All measurements on V100-SXM2-16\,GB.}
\label{tab:width}
\begin{tabular}{rrrrrrl}
\hline
$n_h$ & $d$ & Params\,(MB) & Optim\,(MB) & $C$\,(MB) & $C_\text{EMA}$\,(MB) & $a$\,(MB/bs$\cdot n$) \\
\hline
  2 &   32 &   0.5 &   0.9 &   21 &   22 & 0.0501 \\
  4 &   64 &   1.6 &   3.2 &   26 &   28 & 0.0710 \\
  8 &  128 &   6.1 &  12.3 &   40 &   46 & 0.1129 \\
 16 &  256 &  24.2 &  48.4 &  103 &  127 & 0.1939 \\
 32 &  512 &  96.3 & 192.5 &  329 &  426 & 0.3583 \\
 64 & 1024 & 384.5 & 768.9 & 1222 & 1607 & 0.6921 \\
\hline
\end{tabular}
\end{table}

\begin{figure}[h]
  \centering
  % \includegraphics[width=\textwidth]{vram_plots.pdf}
  \caption{%
    \textbf{VRAM scaling} (six panels).
    \textit{Top row:} peak VRAM vs.\ batch size (log--log) and vs.\
    $n$ particles (linear) at $n_h=16$; lines separate cleanly,
    confirming the bilinear model in Eq.~\eqref{eq:vram}.
    \textit{Middle row:} maximum $n$ before OOM at 16\,GB and 32\,GB
    as a function of batch size, for all six widths.
    \textit{Bottom row:} fixed memory components vs.\ $n_h$ (log--log,
    slope $\approx 2$ confirming $n_h^2$ scaling) and fitted
    activation coefficient $a$ vs.\ $n_h$ (log--log, slope $\approx 1$
    confirming linear scaling).}
  \label{fig:vram}
\end{figure}

\subsubsection{Maximum model width given batch size and particle count}

Table~\ref{tab:maxnh} gives the largest $n_h$ that fits in memory
for each $(\mathrm{bs}, n)$ pair, derived from Eq.~\eqref{eq:vram}
and cross-checked against measured peaks.  Entries in parentheses are
for 32\,GB GPUs; unparenthesised entries fit on 16\,GB.

\begin{table}[h]
\centering
\caption{%
  Maximum $n_h$ (attention heads) that fits within 16\,GB (bold) and
  32\,GB (parenthesised) for each $(\mathrm{bs}, n)$.
  OOM entries are blank.
  Current standard configuration: $n_h=16$, $\mathrm{bs}=512$, $n\leq 6$.}
\label{tab:maxnh}
\begin{tabular}{r|rrrrrr}
\hline
      & \multicolumn{6}{c}{$n$ particles per event} \\
bs    & 2    & 4    & 6    & 8    & 10   & 12   \\
\hline
  256 & \textbf{64}\,(64) & \textbf{64}\,(64) & \textbf{64}\,(64) & \textbf{64}\,(64) & \textbf{64}\,(64) & \textbf{64}\,(64) \\
  512 & \textbf{64}\,(64) & \textbf{64}\,(64) & \textbf{64}\,(64) & \textbf{64}\,(64) & \textbf{64}\,(64) & \textbf{64}\,(64) \\
 1024 & \textbf{64}\,(64) & \textbf{64}\,(64) & \textbf{64}\,(64) & \textbf{64}\,(64) & \textbf{64}\,(64) & \textbf{64}\,(64) \\
 2048 & \textbf{64}\,(64) & \textbf{64}\,(64) & \textbf{64}\,(64) & \textbf{32}\,(64) & \textbf{32}\,(64) & \textbf{32}\,(64) \\
 4096 & \textbf{64}\,(64) & \textbf{32}\,(64) & \textbf{32}\,(64) & \textbf{32}\,(64) & \textbf{16}\,(32) & \textbf{16}\,(32) \\
 8192 & \textbf{32}\,(64) & \textbf{16}\,(32) & \textbf{16}\,(32) & \textbf{8}\,(32)  & \textbf{8}\,(16)  & \textbf{8}\,(16)  \\
\hline
\end{tabular}
\end{table}

Key takeaways:
\begin{itemize}
  \item Below $\mathrm{bs}=2048$, activation memory is small enough
    that the widest model ($n_h=64$) fits for all tested $n\leq 12$
    on 16\,GB.
  \item At $\mathrm{bs}=4096$ (a practical training batch size) one
    can run $n_h=32$ for $n\leq 8$, or $n_h=64$ for $n\leq 2$, on
    16\,GB.
  \item Moving to 32\,GB GPUs (gpu\_p2l) unlocks one width tier at
    every $(bs, n)$: e.g.\ $\mathrm{bs}=8192$, $n=8$ goes from
    $n_h=8$ to $n_h=32$.
\end{itemize}

\subsection{Iteration Time}

\subsubsection{Timing methodology}

Five warmup steps are discarded; 20 steps are timed with
\texttt{torch.cuda.synchronize()} bracketing each step.
Mean and standard deviation are reported; standard deviations are
below $2\,\mathrm{ms}$ in all non-OOM cases.

\subsubsection{Step time}

\begin{table}[h]
\centering
\caption{%
  Step time (ms) at $\mathrm{bs}=2048$ for all measured $(n_h, n)$
  combinations.  OOM cells are blank.
  Throughput $=\mathrm{bs}/t\,[\mathrm{k\,samples/s}]$ in
  parentheses.}
\label{tab:time}
\begin{tabular}{r|rrrrrr}
\hline
     & \multicolumn{5}{c}{$n$ particles per event} \\
$n_h$ & 2 & 4 & 6 & 8 & 10 \\
\hline
 4  & 38\,(53.5) & 51\,(40.4) & 62\,(33.1) & 79\,(26.0) & 93\,(22.1) \\
 8  & 58\,(35.4) & 86\,(23.7) &115\,(17.8) &144\,(14.2) &175\,(11.7) \\
16  &107\,(19.1) &170\,(12.0) &233\,(8.8)  &295\,(6.9)  &360\,(5.7)  \\
32  &226\,(9.0)  &373\,(5.5)  &517\,(4.0)  &664\,(3.1)  &825\,(2.5)  \\
64  &526\,(3.9)  &925\,(2.2)  &1302\,(1.6) &     ---    &     ---    \\
\hline
\end{tabular}
\end{table}

\begin{figure}[h]
  \centering
  % \includegraphics[width=\textwidth]{iteration_time_plots.pdf}
  \caption{%
    \textbf{Iteration time and throughput} (six panels).
    \textit{Top row:} step time (log--log) and throughput vs.\ batch
    size at $n_h=16$, with $n$ as parameter: throughput saturates
    above $\mathrm{bs}\approx 2048$, reaching $\sim 9\,\mathrm{k/s}$
    at $n=6$.
    \textit{Middle row:} same with $n_h$ as parameter at $n=6$:
    wider models are slower but reach lower saturated throughput.
    \textit{Bottom left:} step time vs.\ $n$ at $n_h=16$ is linear
    above $\mathrm{bs}\approx 1024$, confirming MLP dominance
    (FLOPs $\propto n\cdot d^2$).
    \textit{Bottom right:} step time vs.\ $n_h$ at $n=6$ in log--log;
    slope $\approx 1.1$, slightly super-linear at the largest widths.}
  \label{fig:time}
\end{figure}

\subsubsection{Scaling relationships}

\paragraph{Scaling with $n$ (fixed $n_h$, large bs):}
Step time grows linearly with $n$ once $\mathrm{bs}$ is large
enough to saturate the GPU.
This follows from the MLP dominating compute:
$\mathrm{FLOPs}_\text{MLP} \propto n\cdot d^2$, whereas
$\mathrm{FLOPs}_\text{attn} \propto n^2 \cdot d$,
so $\mathrm{FLOPs}_\text{MLP}/\mathrm{FLOPs}_\text{attn}
= d/n = 16\,n_h/n \gg 1$ for $n \leq 12$.
At $n_h=16$, $\mathrm{bs}=2048$:
\[
  \Delta t / \Delta n \approx 31\,\text{ms per extra particle}.
\]

\paragraph{Scaling with $n_h$ (fixed $n$, large bs):}
At $n=6$, $\mathrm{bs}=2048$, doubling $n_h$ increases the step
time by a factor of roughly $2.0$--$2.5$ (slightly super-linear):
\[
  t \propto n_h^{1.1},
\]
consistent with MLP FLOPs $\propto d^2 \propto n_h^2$ partially
offset by increased parallelism at wider widths saturating memory
bandwidth.

\paragraph{Throughput plateau:}
The GPU saturates above $\mathrm{bs}\approx 2048$ for all widths
tested.  Saturated throughput (k\,samples/s) at $n=6$:
$n_h=4$: 35, $n_h=8$: 19, $n_h=16$: 9, $n_h=32$: 4,
$n_h=64$: 1.6.
Throughput scales roughly as $1/n_h$.


% ============================================================
% Behaviour near amplitude divergences & extrapolation
% Research direction opened 2026-07-09 (collaborator-driven).
% Tooling under analysis/divergences/ ; figures in that dir's figs/.
% ============================================================

\section{Behaviour Near Amplitude Divergences and Extrapolation}

So far the model has been judged mostly on \emph{aggregate} accuracy: whether one
Lorentz-equivariant foundation model can \emph{learn} many processes jointly and
\emph{transfer} (fine-tune) to new processes/orders at good pooled validation loss.
Collaborators asked a sharper, physics-driven question: how does the model behave on
the \textbf{divergence structure} of the amplitudes --- the regions where
$|\mathcal{M}|^2$ grows large or varies fastest (soft/collinear limits, $s$-channel
resonances, production thresholds, Sudakov/IR logarithms)?  Does accuracy
\textbf{degrade} there, and can the model \textbf{extrapolate} into those regions
(sparsely-sampled tails and kinematics outside the training support)?

This matters because the peaked / near-divergent regions dominate cross sections and
are the hardest for a surrogate: a foundation model is only useful if it is reliable
\emph{where the physics lives}, not merely on average.  An aggregate
\texttt{val\_loss\_no\_reg} can look excellent while a thin, high-$|\mathcal{M}|^2$
tail is mispredicted.  The evaluation is therefore reframed from a single pooled
number to a \emph{phase-space-local} picture: error as a function of where one sits in
phase space and of how large / how rare the amplitude is.

\subsection{Concrete questions}
\begin{enumerate}
  \item \textbf{Degradation map.} Bin the model error in $\log|\mathcal{M}|^2$ over
    the physical phase space; is the worst error co-located with the divergences
    (low $\sqrt{s}$ / resonance / forward--backward peaks / threshold)?
  \item \textbf{Error vs.\ magnitude.} Does $|\Delta\log|\mathcal{M}|^2|$ rise with the
    $|\mathcal{M}|^2$ percentile (systematically worse on the largest, most physical
    events)?
  \item \textbf{Interpolation vs.\ extrapolation.} Separate ``sparse but in-support''
    from ``outside training support'': fine-tune on a restricted kinematic window
    (e.g.\ $|\cos\theta^*|<0.8$, or $\sqrt{s}$ below a cut) and measure error in the
    held-out window beyond the cut --- true extrapolation, not just rare bins.
  \item \textbf{Specifics.} Near-threshold $e^+e^-\to t\bar t$ ($\sqrt{s}\to 2m_t$),
    the $Z$ pole for $e^+e^-\to u\bar u$, and the genuine soft/collinear IR of the
    real-emission channels ($u\bar u g$, $u\bar u gg$).
\end{enumerate}

\subsection{Fine-tuned NLO virtual: $e^+e^-\to u\bar u$ and $e^+e^-\to t\bar t$}

First deliverable: for the \emph{best} fine-tuned models of the two NLO-virtual
processes (\texttt{ee\_uu\_nlo\_virt\_e4}, \texttt{ee\_ttbar\_nlo\_virt\_e4}; both
$2\to2$, $s$-channel $\gamma^*/Z$), visualise true and predicted
$\log|\mathcal{M}|^2$ across the phase space and localise the error.  Best fine-tunes
(lowest \texttt{val\_loss} over all \texttt{sweeps/*/results/*.json}, both fine-tuned
from \texttt{pretrain\_full\_nh8/trial\_0271}, $D=100$k, $40$k steps):
\texttt{finetune\_scaling\_virt\_002\_eeuunlovirte4\_t40000/trial\_0154} (val
$4.12\mathrm{e}{-}8$, test $9.98\mathrm{e}{-}8$) and the corresponding
\texttt{eettbarnlovirte4} trial (val $2.65\mathrm{e}{-}9$, test $1.25\mathrm{e}{-}7$).

Phase-space coordinates: $\sqrt{s}$ (invariant mass of the initial pair) and the
c.o.m.\ scattering angle $\cos\theta^*$ between the incoming $e^-$ and the outgoing
fermion.  Both are Lorentz invariants, so they are recovered exactly regardless of the
input Lorentz augmentation.

\paragraph{Faithfulness gotcha (record this).}
The extracted per-split MSE must match each run's own end-of-training log before an
error map can be trusted.  The fine-tune runs (2026-06-17) \emph{predate} the
name-based \texttt{amp\_orders} resolver (commit \texttt{accef58}, 2026-06-26): the
stale 8-entry \texttt{amp\_orders} fed \texttt{order\_labels}$=[0,0]$ at training
time, but today's resolver derives $[1,0]$ from ``nlo\_virt'' in the dataset name.
Feeding the wrong constant degraded \emph{every} prediction by $20$--$20000\times$
(even memorised train events).  \texttt{extract\_preds.py} pins the training-time
value; with it, the reproduction matches the logs to 4--5 significant figures.

\paragraph{Result.}  Over $100$k events each, the model reproduces
$\log|\mathcal{M}|^2$ across the \emph{whole} $2\to2$ phase space: RMS
$\Delta\log|\mathcal{M}|^2 \approx 1.7\mathrm{e}{-}4$ ($u\bar u$),
$9.4\mathrm{e}{-}5$ ($t\bar t$).  The near-divergent structure --- the $1/s$ growth
toward low $\sqrt{s}$, the $t\bar t$ threshold, the forward/backward angular peaks ---
is tracked \emph{through} the sharp rises; the largest (still $\sim\!10^{-4}$) errors
sit at the phase-space boundaries ($\cos\theta^*\to\pm1$, $\sqrt{s}$ extremes) and
sparse corners, not at the physical peaks.  Caveat: these virtual $2\to2$ amplitudes
have only \emph{mild} ``divergences'' (finite $1/s$ growth, threshold, angular peaks),
no explicit soft/collinear singularity.

\subsection{Jointly-pretrained $2\to2$ (zero-shot, not fine-tuned)}

The same phase-space view, now for processes the model learned \emph{jointly} in the
8-process pretrain \texttt{pretrain\_full\_nh8/trial\_0271}
($ee\to$ WWZ, WW, $t\bar t$, $u\bar u g$, $u\bar u gg$, $\gamma\gamma$,
$\gamma\gamma\gamma$, $u\bar u$).  Three well-learned $2\to2$ processes with distinct
structures were chosen:

\begin{table}[h]
\centering
\caption{Well-learned jointly-pretrained $2\to2$ processes (zero-shot).}
\begin{tabular}{lll}
\hline
process & pretrain test MSE & divergence structure \\
\hline
$e^+e^-\to\gamma\gamma$ & $9.0\mathrm{e}{-}10$ & $t/u$-channel \textbf{collinear} peaks
  at $\cos\theta^*\to\pm1$ ($\sim1/\sin^2\theta$), $\sqrt{s}$-independent \\
$e^+e^-\to W^+W^-$ & $2.5\mathrm{e}{-}8$ & $t$-channel $\nu$ exchange,
  \textbf{forward} peak growing with $\sqrt{s}$ \\
$e^+e^-\to u\bar u$ & $6.1\mathrm{e}{-}9$ & $s$-channel $\sim(1+\cos^2\theta)$,
  no angular divergence (control) \\
\hline
\end{tabular}
\end{table}

The joint model reproduces all three at RMS $\Delta\log|\mathcal{M}|^2 \approx
3.5\mathrm{e}{-}4$ ($\gamma\gamma$), $4.9\mathrm{e}{-}4$ ($WW$), $4.5\mathrm{e}{-}4$
($u\bar u$).  The collinear divergence of $\gamma\gamma$ (with $\log|\mathcal{M}|^2$
rising to $\approx10$ at $\cos\theta^*\to\pm1$) and the forward peak of $WW$ are
tracked \emph{through} the singular region.  So the ``no degradation at the
divergence, small error only at the phase-space boundary'' behaviour is a property of
the pretrained foundation model, not just of fine-tuning.

\paragraph{Collinear-resolved view.}  Mean-per-bin maps/projections \emph{hide} the
divergence: for $\gamma\gamma$ only $0.003\%$ of events exceed $\log|\mathcal{M}|^2=6$,
all at $|\cos\theta^*|>0.9999$ (peak $9.86$), so a linear-$\cos\theta^*$ bin averages
$\log|\mathcal{M}|^2$ down to $\sim+0.35$ in its outer bin.  Re-plotting against the
signed collinear coordinate
\[
  \eta \;=\; \mathrm{sign}(\cos\theta^*)\,
             \log_{10}\!\frac{1}{1-|\cos\theta^*|}
  \qquad(\eta=0\text{ central},\ \eta\to\pm6\text{ collinear})
\]
stretches both beam directions and turns the divergence into a visible linear ramp
($\log|\mathcal{M}|^2 \sim -\ln(1-|\cos\theta^*|)$).  Plotting mean \emph{and} max per
$\eta$-bin: $\gamma\gamma$ is a symmetric V ramping to $\approx+8$ at both edges; $WW$
an asymmetric forward-only ramp; $u\bar u$ flat (control; the mild forward tilt is the
real $\gamma/Z$ forward--backward asymmetry).  The model tracks truth up the ramps;
error grows only in the most collinear, sparsest bins.

\subsection{Real IR channels: $e^+e^-\to u\bar u g$ ($2\to3$) and
            $e^+e^-\to u\bar u gg$ ($2\to4$)}

The genuine test: gluon-emission channels with real \textbf{soft} ($E_g\to0$) and
\textbf{collinear} ($g\parallel$ quark) singularities, from the same joint pretrain
(zero-shot).  IR observables are built as Lorentz invariants from the raw momenta:
the gluon energy fraction $x_g = 2(p_g\!\cdot\!Q)/s$ (softest $x_{g,\min}$), and the
master resolution variable $y_{\min}=\min_{ij}(p_i+p_j)^2/s$ over gluon-involving
coloured pairs, which $\to0$ in both the soft and collinear limits.

\begin{itemize}
  \item \textbf{$e^+e^-\to u\bar u g$} (pretrain MSE $4.4\mathrm{e}{-}6$): the model
    reproduces the soft ($\sim1/x_g^2$) and soft$+$collinear ($\sim1/y_{\min}$) ramps
    and the \emph{full bremsstrahlung Dalitz} (collinear walls at $x_q\to1$,
    $x_{\bar q}\to1$ meeting at the soft $(1,1)$ corner) --- truth $\approx$ model
    across $\sim\!30$ orders of magnitude in $|\mathcal{M}|^2$, RMS
    $\Delta\log|\mathcal{M}|^2 = 0.007$.  No IR degradation.
  \item \textbf{$e^+e^-\to u\bar u gg$} (pretrain MSE $2.3\mathrm{e}{-}4$, the hardest
    process in the set): the model still tracks the \emph{mean} IR ramps into the deep
    IR, but this is the \textbf{first clear case of divergence-region degradation} ---
    the per-event error grows monotonically into the IR, $\langle
    |\Delta\log|\mathcal{M}|^2|\rangle \approx 0.03$ at $y_{\min}\!\sim\!0.1
    \to 0.27$ at $y_{\min}\!\sim\!10^{-4}$, and the pred-vs-true band widens (RMS
    $0.085$).  The error maps localise it to the deep soft/collinear corner.
\end{itemize}

\paragraph{Bottom line.}  Across all $2\to2$ (including the $\gamma\gamma$ collinear
divergence) and $e^+e^-\to u\bar u g$, accuracy does \emph{not} degrade at the
divergences --- only at phase-space boundaries / sparse bins.  Degradation
\emph{into} the singular region appears only for the hardest, highest-multiplicity
process ($u\bar u gg$) in its deepest soft/collinear limits.

\subsection{Tooling and next steps}

Everything lives under \texttt{analysis/divergences/}: \texttt{extract\_preds.py}
(fine-tune runs), \texttt{extract\_pretrain.py} (jointly-pretrained $2\to2$),
\texttt{extract\_ir.py} (gluon channels) --- each reloads a run's exact config /
preprocessing / $\mu$P model, loads the checkpoint, forwards on GPU (\texttt{sbatch};
the xformers block-diagonal attention is CUDA-only), and recovers exact physical
kinematics by inverting the deterministic \texttt{default\_rng(42)} event shuffle back
to the raw rows.  They are preprocessing-scope-aware (GLOBAL vs.\ PER-DATASET amp
stats, \texttt{len(prepd\_mean)>1}) and guard alignment/stats with a per-event
$\mathrm{preprocess}(\mathrm{raw\_amp})=$ stored-truth check (max$|\Delta|=0$).
\texttt{make\_plots.py} draws the maps / projections / hexbin \emph{plus} the
collinear-resolved row; \texttt{make\_ir.py} the IR ramps + Dalitz.

Next: (i) an error-vs-$|\mathcal{M}|^2$-percentile curve (question 2); (ii) the
restricted-window fine-tune for a clean extrapolation number (question 3); (iii)
whether a short fine-tune (or more pretrain weight on $u\bar u gg$) recovers the deep
soft/collinear limit where degradation was seen.
